\scanpage{24}
\mspage{- 2 -}
\margin{n°355}

d'où $\operatorname{Ass}_B(F/\underline pF)\subset\operatorname{Ass}_B(E\otimes_A F)$, ce qui démontre l'inclusion $(\ast)$.

On va maintenant supposer que $A$ est \uline{noethérien}, et démontrer l'inclusion
opposée. On procèdera par étapes :

(i) $E$ \uline{est un module de type fini et} $\operatorname{Ass}_A(E)$ \uline{est réduit à un élément} $\underline p$.

D'après le th. 1 du §1, il existe une suite de composition $(E_i)_{0\leq i\leq n}$ de
$E$ telle que $E_i/E_{i+1}$ soit isomorphe à $A/\underline p_i$, où $\underline p_i$ est un idéal premier de $A$.
D'après le th.2 du §1, tous les $\underline p_i$ contiennent $\underline p$. Du fait que $F$ est $A$-plat,
les $E_i\otimes_A F$ forment une suite de composition de $E\otimes_A F$, et
$(E_i\otimes_A F)/(E_{i+1}\otimes_A F)$ s'identifie à $A/\underline p_i\otimes_A F=F/\underline p_iF$. En appliquant la
prop.3, on en déduit :
\[
\operatorname{Ass}_B(E\otimes_A F)\subset\bigcup\operatorname{Ass}_B(F/\underline p_iF).
\]

Soit donc $\underline P\in\operatorname{Ass}_B(E\otimes_A F)\cap\operatorname{Ass}_B(F/\underline p_iF)$. Si $\underline p_i\ne\underline p$, il existe $a\in\underline p_i$,
$a\notin\underline p$. D'après le cor.2 à la prop.2 du §1, l'homothétie de
rapport $a$ est injective dans $E$, donc celle de rapport $\pi(a)$ est injective
dans $E\otimes_A F$ ; mais d'autre part, $\pi(a)$ annule $F/\underline p_iF$, d'où $\pi(a)\in\underline P$, contradiction.
Ainsi $\underline p_i$ est nécessairement égal à $\underline p$, ce qui démontre la proposition dans le cas considéré.

(ii) $E$ \uline{est un module de type fini}.

Soient $\underline p_i$ les éléments de $\operatorname{Ass}_A(E)$, et soit $0=\bigcap Q_i$ une décomposition
primaire réduite correspondante (cf. §2, n°2). Le module $E$ se plonge dans
la somme directe des $E_i=E/Q_i$, et $E\otimes_A F$ se plonge dans la somme directe
des $E_i\otimes_A F$. On en déduit :
\[
\operatorname{Ass}_B(E\otimes_A F)\subset\bigcup\operatorname{Ass}_B(E_i\otimes_A F).
\]
D'après (i), on a $\operatorname{Ass}_B(E_i\otimes_A F)=\operatorname{Ass}_B(F/\underline p_iF)$, d'où la proposition dans
ce cas.

(iii) \uline{Cas général}.

\margin{23}
